% TOP_PROBLEM: {"id": 1100214, "title": "Questions in Geometric Group Theory — Q 2.14", "category_id": 17, "category": {"id": 17, "name": "group_theory", "display_name": "Group Theory", "description": "Problems about groups, group actions, representations, and related algebraic structures.", "slug": "group-theory", "order_index": 17, "created_at": "2026-07-31T15:26:25.671Z"}, "status": "open", "problem_number": "AMR-010-0214", "set_id": 13}
% FIRST_POSTED: 2026-10-01
% ENTRY_KIND: statement_only
% UnsolvedMath ID 1100214; source catalogue code AMR-010-0214
% !TeX program = lualatex
% Definition and sources only; this file is not an LLM solution attempt.
\documentclass[11pt,a4paper]{article}
\usepackage[margin=25mm]{geometry}
\usepackage{fontspec}
\setmainfont{lmroman10-regular.otf}[BoldFont=lmroman10-bold.otf,
 ItalicFont=lmroman10-italic.otf,BoldItalicFont=lmroman10-bolditalic.otf]
\usepackage{amsmath,amssymb,mathtools}
\usepackage{unicode-math}
\setmathfont{latinmodern-math.otf}
\AtBeginDocument{\let\setminus\smallsetminus}
\usepackage{xurl,hyperref}
\hypersetup{colorlinks=true,urlcolor=blue}
\setcounter{secnumdepth}{0}
\setlength{\parindent}{0pt}
\setlength{\parskip}{5pt}
\setlength{\emergencystretch}{3em}
\begin{document}
\section*{Questions in Geometric Group Theory — Q 2.14}\label{1100214}
{\small UnsolvedMath ID 1100214 · AMR-010-0214 · Group Theory · AMR Open Problem Lists}
\subsection{Definitions and mathematical statement}
Fix a finite set $S$ and symmetric labels $m_{st}\in\{2,3,\ldots\}\cup\{\infty\}$ for its distinct pairs. Write $[s,t]_m$ for the alternating word of length $m$ starting with $s$, so $[s,t]_3=sts$. The associated Artin group has presentation
\[
 A(M)=\langle S\mid [s,t]_{m_{st}}=[t,s]_{m_{st}}
                      \text{ for }m_{st}<\infty\rangle.
\]
An infinite label contributes no relation, and the generators have no involution relations. The input is the finite matrix of labels.

Determine when two such groups $A(M)$ and $A(N)$ are isomorphic as abstract groups. The task is a complete characterization in terms of the defining data, with an effective recognition criterion if possible. There is no requirement that an isomorphism respect standard generators, their conjugacy classes, or a chosen diagram decomposition.

Consequently, simply testing whether the labeled diagrams are isomorphic does not itself state the classification problem. Conversely, any proposed moves on diagrams must be proved sufficient and necessary for abstract group isomorphism if they are to give a full answer. The source question includes arbitrary finite Artin matrices, not just right-angled groups, matrices with large finite labels, or groups of spherical type. Here spherical type means that adding $s^2=1$ for all standard generators gives a finite Coxeter group. The possible rank and the number of infinite labels are unrestricted finite parameters.
\subsection{Short English statement}
Classify all finitely generated Artin groups up to abstract group isomorphism from their defining matrices.
\subsection{Sources}
\begin{itemize}
\item[S1] ULAM.AI and the UnsolvedMath Contributors, supplied problems.json, record 1100214 (AMR-010-0214), AMR Open Problem Lists. Catalogue identity and source transcription; CC BY 4.0.
\url{https://huggingface.co/datasets/ulamai/UnsolvedMath}.
\item[S2] Bestvina - Questions in Geometric Group Theory (pdf) (2004), Question 2.14 (PDF page 9). Original source indexed by AMR-010-0214.
\url{https://www.math.utah.edu/~bestvina/eprints/questions-updated.pdf}.
\end{itemize}
\end{document}
